\section{Applications by field}\label{sec:applications}
The theories reviewed here serve five broad application areas. For each, the relevant regime of parameter space is stated, using the characteristic numbers derived above.

\subsection{Laser ultrasonics and short-pulse heating}
Laser-generated ultrasound uses the thermal stress of pulsed heating to generate elastic waves \citep{Scruby1990}. For nanosecond pulses on bulk metals $\omega\tau_q\ll1$, the classical coupled theory with surface heating suffices, and the thermal diffusion length over the pulse, $\sqrt{\alpha t_p}$, sets the source depth. For femtosecond--picosecond pulses the heat-carrier relaxation, electron--phonon coupling in metals, falls on time scales where hyperbolic or two-temperature models matter, and thin-film heating is modeled with lagging, fractional or memory-dependent laws; bilayer analyses under pulsed heating combine them \citep{Li2019}. The benchmark thermal shock of Section~\ref{sec:transient} shows what to look for: the stress peak at the elastic front and its sensitivity to the heat law, here about 20\,\% between Fourier and Lord--Shulman for the parameters used.

\subsection{Micro- and nano-electromechanical resonators}
Thermoelastic damping is a leading loss mechanism in flexural micro-resonators. The classical Zener relaxation (Section~\ref{sec:ted}) captures the dependence on thickness and frequency, and the extension to beams and plates by Lifshitz and Roukes \citep{Lifshitz2000} is standard. The corrections reviewed here enter in two ways. In ultra-thin structures at very high frequency, lagging and non-Fourier effects change the damping; examples are DPL-based thermoelastic damping in nonlocal nanoplates \citep{Yang2021}, couple-stress micro-plates \citep{Kumar2017} and rotating nano-beams with spatiotemporal nonlocality \citep{Sur2025}. At small sizes, nonlocal and gradient elasticity also shift the vibration frequencies on which the damping depends. Two cautions follow from the analysis presented here: the form of the nonlocal thermoelastic constitutive law (form A or B, Section~\ref{sec:size}) changes the predicted damping near the elastic cut-off, and the lag time enters only through $\omega\tau$, so unless the frequency is within one or two decades of $1/\tau_q$, the classical result is adequate.

\subsection{Piezoelectric devices, sensors and semiconductors}
Thermopiezoelectric coupling affects bulk and surface acoustic wave devices, energy harvesters and sensors, and photothermal effects in piezoelectric semiconductors add carrier dynamics \citep{Gupta2023,Gupta2024a,Gupta2026HeatMap}. In the quasi-static, charge-free regime the electric field acts on thermoelastic waves as the exact renormalization \eqref{eq:renorm}, and designs that sense temperature through the pyroelectric effect depend on the sign convention discussed there. Interface problems (substrate--film, electrode--crystal) are the subject of the energy-ratio literature \citep{Barak2023IJP,Barak2023JTS,Gupta2023Energy,Kumar2024,Kumar2024a}, where Section~\ref{sec:interfaces} suggests the checks to apply.

\subsection{Biological tissue and heterogeneous media}
Heat transfer in biological tissue is classically described by the bioheat equation of Pennes \citep{Pennes1948}, which adds blood perfusion to the Fourier equation. Non-Fourier models have been applied to tissue, on the strength of reported large relaxation times in nonhomogeneous materials \citep{Kaminski1990,Mitra1995}, though the interpretation of those measurements is debated; DPL and fractional models have been fitted to such data \citep{vanDerMerwe2021,Yu2024}. Shock-induced nonlocal coupled thermoelasticity analysis with energy dissipation in a MEMS/NEMS beam resonator was presented by \citep{Hosseini2017}. Fractional order thermo-viscoelastic theory of biological tissue with dual phase lag heat conduction model was developed by \citep{Li2021}; bio-heat transfer analysis based on fractional derivative and memory-dependent derivative heat conduction models was performed by \citep{Zhang2021}. For water-like media $\omega_*=\rho cc_1^2/k\approx10^{13}$\,rad/s (with $\rho\approx10^3$\,kg/m$^3$, $c\approx4\times10^3$\,J/kg\,K, $c_1\approx1.5\times10^3$\,m/s, $k\approx0.6$\,W/m\,K) lies far above any frequency of interest, and the thermoelastic coupling is weak, so tissue work concerns the thermal field rather than thermoelastic waves.

Effect of magnetic field and gravity on 2D fiber-reinforced medium under fractional order theory of thermoelasticity was analyzed by \citep{Ma2020}.

\subsection{Geophysical and porous media}
Porous and double-porous thermoelastic models with two temperatures are applied to wave propagation in saturated and partially saturated solids \citep{Singh2020,Kumar2022,Kumar2025}, where additional slow waves carry fluid and heat and attenuate strongly. The validity of generalized heat laws in such media, whose time scales are set by pore-scale processes, is largely unexamined.

\begin{figure}[t!]\centering\small
\input{figs_tikz/fig11_applications.tex}
\caption{Application regime map: five areas ordered by the dimensionless numbers $\omega\tau$ (lag relevance) and $\omega/\omega_*$ (coupling relevance); Section~\ref{sec:synthesis} gives the quantitative criteria.}\label{fig:applications-geo}
\end{figure}

\subsection{Summary}
In all five areas the single most informative dimensionless number is $\omega\tau$ (frequency times the relevant lag), and the second is $\omega/\omega_*$. The regime map in Section~\ref{sec:synthesis} makes this quantitative.

\begin{summarybox}[Summary]
Applications divide by regime: laser ultrasonics (classical adequate for ns pulses; lagged/fractional for fs--ps), MEMS/NEMS resonators (Zener damping with nonlocal/lag corrections near cut-off), piezoelectric devices (exact stiffening renormalization in charge-free regime), biological tissue (weak coupling, interest in thermal field), geophysical porous media (slow diffusive waves, largely unexamined). The key numbers are $\omega\tau$ and $\omega/\omega_*$.
\end{summarybox}
